\documentclass[10pt]{amsart}
\usepackage[a4paper,margin=23mm]{geometry}
\usepackage{amsmath,amssymb,amsthm,xfrac,hyperref,xspace,graphicx}
\usepackage[english]{babel}
\usepackage[scr=rsfs,cal=euler]{mathalfa}
\newcommand{\resp}{resp.}
\newcommand{\cf}{cf.}
\newcommand{\ie}{i.e.}
\newcommand{\eg}{e.g.}
\newcommand{\loccit}{loc.\ cit.}
\newcommand{\skpt}{\par\smallskip\noindent}
\let\Subsection\subsection
\let\Subsubsection\subsubsection
\emergencystretch=3em
\pagestyle{plain}
\multlinegap0pt
\hyphenation{asso-cia-ted assu-me espe-cial-ly infor-ma-tion inves-ti-gate opti-mal}

\newcommand\mto{\mathchoice{\longmapsto}{\mapsto}{\mapsto}{\mapsto}}
\newcommand{\Babs}[1]{\Bigl|#1\Bigr|}
\newcommand{\Bnorm}[1]{\Bigl\|#1\Bigr\|}
\newcommand{\psfrac}[2]{\sfrac{(#1)}{#2}}
\newcommand{\spfrac}[2]{\sfrac{#1}{(#2)}}
\newcommand{\Psfrac}[2]{(\sfrac{#1}{#2})}
\newcommand{\Ppsfrac}[2]{(\psfrac{#1}{#2})}

\usepackage{pgfplots}
\usepackage{subfigure}
\pgfplotsset{compat=1.15}
\usepackage{dsfont}
\usepackage{mathtools}
\usepackage{tikz}

\newtheorem{theorem}{Theorem}
\newtheorem{corollary}[theorem]{Corollary}
\newtheorem{lemma}[theorem]{Lemma}
\newtheorem{proposition}[theorem]{Proposition}
\newtheorem{example}[theorem]{Example}

\theoremstyle{definition}
\newtheorem{definition}[theorem]{Definition}
\newtheorem{definition-proposition}[theorem]{Definition-Proposition}
\newtheorem*{notation}{Notation}
\newtheorem{remark}[theorem]{Remark}

\newcommand{\N} {\mathbb{N}}
\newcommand{\Z} {\mathbb{Z}}
\newcommand{\C} {\mathbb{C}}
\newcommand{\R} {\mathbb{R}}

\newcommand{\yp}[1]{{#1}}

\DeclarePairedDelimiter\abs{\lvert}{\rvert}
\DeclarePairedDelimiter\norm{\lVert}{\rVert}
\newcommand{\tnorm}[1]{{\left\vert\kern-0.25ex\left\vert\kern-0.25ex\left\vert #1
\right\vert\kern-0.25ex\right\vert\kern-0.25ex\right\vert}}

\makeatletter
\let\oldabs\abs
\def\abs{\@ifstar{\oldabs}{\oldabs*}}
\let\oldnorm\norm
\def\norm{\@ifstar{\oldnorm}{\oldnorm*}}
\makeatother

\newcommand*\diff{\mathop{}\!\mathrm{d}}
\newcommand*\Diff[1]{\mathop{}\!\mathrm{d^#1}}
\newcommand{\Tr}{\operatorname{Tr}}
\newcommand{\sinc}{\operatorname{sinc}}

\begin{document}
\title{Strict positive control threshold for minimal-time uniform Landau–Lifshitz magnetization reversal in ellipsoids with a simple easy axis}
\author{Raphaël Côte \and Clémentine Courtès \and Guillaume Ferrière \and Yannick Privat}
\maketitle
\section{Source equation and control setting}
with $\alpha>0$, a constant (in time and space) damping coefficient which is characterized by the material. We refer for instance to \cite{hubert2008magnetic,brown1963micromagnetics} for additional explanations.
As a consequence, the Landau-Lifshitz equation with a time-dependent external magnetic field $h_{\rm ext}$ reads as the ordinary differential system
\begin{equation}\label{LL:ODE}
\begin{cases}
\dot{m} = \alpha \left(h_0(m)-(h_0(m)\cdot m) m\right)-m\wedge h_0(m) & \text{in }(0,T), \\
m(0) =m^0,
\end{cases}
\end{equation}
where the dotted notation $\dot{m}$ stands for the time derivative of $m$, $h_0(m)=-Dm+h_{\rm ext}$, $T>0$, $m(t)\in\mathbb{S}^2\subset\R^3$, $D=\operatorname{diag}([\gamma_1,\gamma_2,\gamma_3])$ is a diagonal matrix with positive coefficients.
Up to a change of basis, we~will also assume without loss of generality that
\begin{equation}\label{hyp:alpha}
0\leq \gamma_1\leq \gamma_2\leq\gamma_3\leq 1.
\end{equation}
\section{Optimal control problem}\label{sec:OCP}
\setcounter{theorem}{1}

The main issue we want to tackle reads
\begin{quote}
\textit{Given a steady-state $\overline{m}$ of \eqref{LL:ODE} in $\mathbb{S}^2$, can we achieve a reversal by solving an optimal control problem, \ie steering the system from $m(0)=\overline{m}$ to $m(T)=-\overline{m}$ while minimizing $T$?}
\end{quote}
In what follows, we~will consider particular stationary states: $\overline{m}=\pm e_1$. It~is proved in appendix~\ref{append:e1AS} that these equilibria are asymptotically stable when $\gamma_1 < \gamma_2$. Therefore, they can be used as magnetic spin orientation for memory storage purposes.
We will denote by $u(\cdot)$ the external (spatially uniform) magnetic field imposed on the system. This is the control variable in this problem. The question is then to ask if it is possible to steer the solution $m_u$ of the system \eqref{LL:ODE} associated to the control $u(\cdot)$ and to the initial data $m_u(0)=e_1$ until $m_u(T)=-e_1$, in~minimal time.

Of course, it is necessary to add physical constraints to this problem: if one imposes no restrictions on the choice of admissible controls, it is likely that the minimal time problem will have a solution, reached by unrealistic controls\footnote{{Indeed, an applied magnetic field in a region cannot be excessively strong, as it is not feasible to create infinitely strong fields due to physical limitations such as material saturation and thermal effects. }}.
For this reason, we~will assume in what follows the constraint
\begin{equation}\label{constr:barU}\tag{$\mathscr{C}_{T,U }$}
|u(t)|\leq U \quad \text{a.e. in }(0,T),
\end{equation}
in order to limit the choice of controls to realistic possibilities.

All in all, the problem we aim at investigating reads as follows.

\Subsubsection*{Minimal time problem}
\begin{quote}
\textit{Let $U >0$ and assume that $m_0=e_1$. The problem reads
\begin{equation}\label{OCP:minT}\tag{$\mathscr{P}_0$}
\mathcal{T}_{U } \coloneqq \inf_{(T,u)\in \mathcal{O}_{U }}T,
\end{equation}
where
\[
\mathcal{O}_{U }=\{(T,u)\mid T\in \R_+, \ u\in L^\infty(0,T)\text{ satisfies }\eqref{constr:barU}\text{ and }m_u(T)=-e_1\},
\]
with $m_u$, the solution to \eqref{LL:ODE} associated to the control function $u(\cdot)$ and the initial datum $e_1$.}
\end{quote}

We will investigate the following issues:
\begin{itemize}
\item Does Problem~\eqref{OCP:minT} have an optimal solution for any value of $U >0$ ?
\item How to characterize all the solutions to this problem and understand their geometric dependence to the parameters $\gamma_i$, $i=1,2,3$?
\end{itemize}

\subsection{Main results}\label{sec:mainRes}

First, the minimization problem is indeed well-posed, meaning that the existence of an optimal solution is equivalent to the existence of a minimal trajectory.

\begin{theorem} \label{th:ex_min}
Let $U > 0$. The following properties are equivalent:
\begin{itemize}
\item[(i)] There exists an optimal pair $(\mathcal{T}_{U}, u) \in \mathcal{O}_{U }$ for Problem~\eqref{OCP:minT}.
\item[(ii)] $\mathcal{T}_{U}$ is finite.
\item[(iii)] $\mathcal{O}_{U }$ is nonempty.
\end{itemize}
\end{theorem}

\begin{theorem} \label{th:ex_U_crit}
Assume $\gamma_1 < \gamma_2$. Then there exists $U_\textnormal{crit} > 0$ such that
\begin{itemize}
\item for all $U \in (0, U_\textnormal{crit}]$, \eqref{OCP:minT} has no solution.
\item for all $U > U_\textnormal{crit}$, \eqref{OCP:minT} has a solution.
\end{itemize}
\end{theorem}

\section{Optimal-trajectory existence}
\Subsection{Proof of Theorem~\ref{th:ex_min}: existence of an optimal trajectory}

Let us assume that~$\mathcal{O}_{U }$ is nonempty. This allows us to consider a minimizing \hbox{sequence} $(T_n,u_n)_{n\in\N}\in\nobreak \mathcal{O}_{U }^\N$, and $m_n \in \mathscr{C}([0,T_n])$ the solution to \eqref{LL:ODE} with field $u_n$.
By~definition, $T_n \to \mathcal{T}_U$ as $n \to \infty$.
In what follows, we~will index similarly a sequence and any subsequence with a slight abuse of notation, for the sake of simplicity.

Let us introduce the functions $\widetilde{u}_n$, $\widetilde{m}_{n}$ defined on $[0,1]$ by
\[
\widetilde{u}_n(s)=u_n(T_ns)\quad \text{and}\quad \widetilde{m}_{n}(s)=m_{n}(T_ns).
\]
Hence, System~\eqref{LL:ODE} rewrites
\begin{equation}\label{LL:ODEbis}
\begin{cases}
\dot{\widetilde{m}}_n = T_n \Bigl(\alpha \bigl(\widetilde{h}_0(\widetilde{m}_n)-(\widetilde{h}_0(\widetilde{m}_n)\cdot \widetilde{m}_n) \widetilde{m}_n\bigr)-\widetilde{m}_n\wedge \widetilde{h}_0(\widetilde{m}_n)\Bigr)& \text{in }(0,1), \\
\widetilde{m}_n(0) =e_1,
\end{cases}
\end{equation}
where $\widetilde{h}_0(\widetilde{m}_n)=-D\widetilde{m}_{n}+\widetilde{u}_n$.

Similarly, since the sequence $(\widetilde u_n)_{n\in\N}$ is bounded in $L^\infty(0,1)$, it converges weakly-star in $L^\infty(0,1)$ up to a subsequence to some element $u^*$ such that $|u^*(\cdot)|\leq U $ a.e. in $[0,1]$ according to the Banach-Alaoglu-Bourbaki theorem. Since both $(\widetilde m_{n})_{n\in \N}$ and $(\widetilde u_n)_{n\in \N}$ are bounded in $L^\infty([0,1])$, we~infer that so is $\dot{\widetilde m}_{n}$ according to \eqref{LL:ODEbis}. Therefore, the sequence $(\widetilde m_{n})_{n\in \N}$ is bounded in $W^{1,\infty}(0,1)$ and hence converges (up~to a subsequence) towards an element $\widetilde{m}^*\in W^{1,\infty}(0,1)$ in $\mathcal{C}^0([0,1])$ according to the Ascoli theorem. In particular, one has necessarily $|\widetilde{m}^*(\cdot)|=1$. Now, let us rewrite \eqref{LL:ODEbis} as the fixed-point equation: for all $s\in [0,1]$,
\[
\widetilde m_{n}(s)=e_1+T_n\int_0^s\Bigl(\alpha \bigl(\widetilde h_0(\widetilde m_{n})-(\widetilde h_0(\widetilde m_{n})\cdot \widetilde m_{n}) \widetilde m_{n}\bigr)-\widetilde m_{n}\wedge \widetilde h_0(\widetilde m_{n})\Bigr)\, d\sigma .
\]
Observe that the right-hand side is linear with respect to $\widetilde h_0(m_{n})$ and that, according to the properties above, $(\widetilde h_0(\widetilde m_{n}))_{n\in \N}$ converges weakly-star to $\widetilde h_0(\widetilde m^*)$ in $L^\infty(0,1)$, where $\widetilde{h}_0(\widetilde{m}^*)=-D\widetilde{m}^*+\widetilde{u}^*$. Letting $n$ tend to $+\infty$ in the equation above, we~obtain:
\[
\forall s\in [0,1], \;\; \widetilde m^*(s)=e_1+\mathcal{T}_U\int_0^s\Bigl(\alpha \bigl(\widetilde h_0(\widetilde m^*)-(\widetilde h_0(\widetilde m^*)\cdot \widetilde m^*) \widetilde m^*\bigr)-\widetilde m^*\wedge \widetilde h_0(\widetilde m^*) \Bigr)\, d\sigma.
\]
Moreover, since $\tilde{m}_n (1) = - e_1$ by construction, the convergence in $\mathcal{C}^0 ([0, 1])$ leads to $\tilde{m}^* (1) = - e_1$. Taking the previous formula with $s=1$, we~get
\begin{equation*}
- 2 e_1 = \mathcal{T}_U\int_0^1\Bigl(\alpha \bigl(\widetilde h_0(\widetilde m^*)-(\widetilde h_0(\widetilde m^*)\cdot \widetilde m^*) \widetilde m^*\bigr)-\widetilde m^*\wedge \widetilde h_0(\widetilde m^*)\Bigr)\, d\sigma.
\end{equation*}
This proves that $\mathcal{T}_U > 0$.
Now, let us introduce $u^*$ as $u^*(t)=\widetilde{u}^*(t/\mathcal{T}_U)$. By~undoing the change of variable above, we~get that $\widetilde{m}^*(t/\mathcal{T}_U)=m_{u^*}(t)$ for a.e. $t\in [0,\mathcal{T}_U]$. Furthermore, $m_{u^*}(0)=e_1$ and $m_{u^*}(\mathcal{T}_U)=-e_1$ since $\widetilde{m}_n(0)=e_1$ and $\widetilde{m}_n(1)=-e_1$. The converse sense is straightforward and the expected conclusion follows.

Finally, observe that the same reasoning can be reproduced whenever $\mathcal{T}_U$ is finite.
\setcounter{theorem}{9}
\subsection{Sufficient and necessary condition for the existence of an admissible trajectory}

Let $(m,u)$ be a solution to \eqref{LL:ODE} on $[0,T]$, and for $t \in [0,T]$, consider the mobile frame $\mathcal{B}(t)=(m(t),\dot e(t),m(t)\wedge \dot e(t))$ where\footnote{Here, $\dot e$ is merely a notation, and not the time derivative of a previously defined vector $e$.} $\dot{e}=\dot{m}/|\dot{m}|$. According to \cite{alouges2009magnetization}, by~observing that
\[
m \perp \dot m, \quad \dot m \perp m \wedge \dot m \quad \text{and}\quad m \perp m \wedge (m\wedge Dm),
\]
one shows easily, by~decomposing $u(t)$ into $\mathcal{B}(t)$ and writing the equation for \hbox{$u - (u \cdot m) m$}, the projection of $u$ on $m^\perp$, that there exists $\lambda\in L^\infty (0,T)$ such that
\begin{equation} \label{eq:u_m}
u = \frac{1}{1+\alpha^2}\,(\alpha \dot m + m \wedge \dot m)+Dm-(Dm\cdot m)m+\lambda m.
\end{equation}
In fact, $\lambda = u\cdot m$. Reciprocally, given any function $m \in W^{1,\infty}([0,T],\mathbb S^2)$, and any function $\lambda \in L^\infty([0,T],\mathbb R)$, if we define $u$ by \eqref{eq:u_m}, then $(m,u)$ is solution to \eqref{LL:ODE}. These considerations can be seen as a consequence of a flatness property of the main system.

Again assuming that $(m,u)$ is admissible trajectory, \ie a solution to \eqref{LL:ODE}, We infer from \eqref{eq:u_m} that $u(t)$ expands as
\begin{equation}\label{eq:uexpand}
u=\lambda m +\Bigl(\frac{\alpha}{1+\alpha^2}|\dot m|+\dot e\cdot Dm \Bigr)\dot e+\Bigl(Dm \cdot (m \wedge \dot e)+\frac{1}{1+\alpha^2}|\dot m|\Bigr)m \wedge \dot e.
\end{equation}
As, a consequence, using that $Dm \cdot (m \wedge \dot e)=\dot e\cdot (Dm \wedge m)$ due the triple product property, we~get
\begin{align*}
|u|^2 &= \lambda^2+\Bigl(\frac{\alpha}{1+\alpha^2}|\dot m|+\dot e\cdot Dm\Bigr)^2+\Bigl(Dm\cdot (m\wedge \dot e)+\frac{|\dot m|}{1+\alpha^2}\Bigr)^2\\
&= \lambda^2+\Bigl(\dot e\cdot Dm+\frac{\alpha|\dot m|}{1+\alpha^2}\Bigr)^2+\Bigl(\dot e\cdot (Dm_u\wedge m)+\frac{|\dot m|}{1+\alpha^2}\Bigr)^2.
\end{align*}
Clearly, this computation and the previous remarks show that, without loss of generality, we~can furthermore assume that an \emph{optimal} trajectory satisfies $\lambda=0$, or equivalently, $u \cdot m =0$: we will do this in the following.

Let us introduce, for a given $T>0$,
\[
\mathscr{V}_T=\{m\in H^1([0,T];\mathbb{S}^2)\mid m(0)=e_1\text{ and }m(T)=-e_1\}.
\]

To investigate the existence of an admissible trajectory, it is then convenient to introduce
\begin{align}\label{ob:aux}
L_\alpha(m) &\coloneqq \Bigl(\dot e\cdot (Dm\wedge m)+\frac{|\dot m|}{1+\alpha^2}\Bigr)^2+\Bigl(\dot e\cdot Dm+\frac{\alpha|\dot m|}{1+\alpha^2}\Bigr)^2\\
\Lambda(T)&\coloneqq \inf_{\lambda\in L^\infty([0,T])}\inf_{m\in \mathscr{V}_T}\sup_{t\in [0,T]}\left(\lambda^2+L_\alpha(m) \right)\nonumber \\
&= \inf_{m\in \mathscr{V}_T}\sup_{t\in [0,T]}L_\alpha(m).
\end{align}

We summarize the above discussion in the form of a lemma.

\begin{lemma}\label{lem:2057}
The existence of an admissible trajectory for Problem~\eqref{OCP:minT} comes to the existence of $m\in \mathscr{V}_T$ such that the function $u$ given by \eqref{eq:uexpand} with $\lambda =0$ satisfies $\Vert u\Vert_{L^\infty([0,T];\mathbb{R}^3)}\leq U $, which is also equivalent to $\Lambda(T)\leq U ^2$. Also, $u$ satisfies $u \cdot m=0$.
\end{lemma}

\subsection{Necessary optimality conditions for Problem~\texorpdfstring{\eqref{OCP:minT}}{P0}}\label{optCond:PMP}
This problem can be solved by using the Pontryagin maximum principle.
\setcounter{theorem}{12}
\section{Proof of Theorem~\ref{th:ex_U_crit}}\label{proof:theo2}
\subsection{Preliminary results}
We first state preliminary results, in~the form of a series of lemmas.

\begin{lemma} \label{lem:cont_flow}
For all $T < \infty$, the map
\begin{align*}
L^\infty (0, T) &\to W^{1, \infty} (0, T) \\
u &\mto m \text{ solution to \eqref{LL:ODE} with } m(0) = e_1
\end{align*}
is continuous and locally Lipschitz.
\end{lemma}

\begin{proof}
Let $u_1, u_2 \in L^\infty (0, T)$ and $m_1, m_2$ the corresponding solution to \eqref{LL:ODE}. Define $\delta m \coloneqq m_1 - m_2$ and $\delta u \coloneqq u_1 - u_2$.
Simple, though tedious, calculations provide that~$\delta m$ satisfies
\begin{multline*}
\frac{d\delta m}{dt} = \alpha \Bigl[ - D \, \delta m + \delta u - ((- D \, \delta m + \delta u) \cdot m_1) m_1 \\[-5pt]
- ((- D m_2 + u_2) \cdot \delta m) m_1 - ((- D m_2 + u_2) \cdot m_2) \delta m \Bigr] \\ - m_1 \wedge (- D \delta m + \delta u) - \delta m \wedge (- D m_2 + u_2).
\end{multline*}
in $(0,T)$. Since $\abs{m_1} = \abs{m_2} = 1$, we~obtain
\begin{equation} \label{eq:est_delta_m_dot}
\Babs{\frac{d\delta m}{dt}} \leq \left((4 \alpha + 2) \norm{D}_2 + (2\alpha+1) \abs{u_2} \right) \abs{\delta m} + (2 \alpha + 1) \abs{\delta u}, \quad \text{in }(0,T),
\end{equation}
where $\norm{\cdot}_2$ stands for the operator norm associated to the euclidean norm $\abs{\cdot}$.
Since $\delta m (0) = 0$, we~have for all $t \in (0, T)$
\begin{multline*}
\abs{\delta m (t)}\leq \int_0^t \Babs{\frac{d\delta m}{dt}} (s) \diff s \leq (2 \alpha + 1) T \norm{\delta u}_{L^\infty} \\
+\int_0^t \left((4 \alpha + 2) \norm{D}_2 + (2\alpha+1) \norm{u_2}_{L^\infty} \right) \abs{\delta m} (s) \diff s,
\end{multline*}
and thus by Gronwall's lemma,
\begin{equation*}
\abs{\delta m (t)} \leq (2 \alpha + 1) T \norm{\delta u}_{L^\infty} \exp{\bigl(((4 \alpha + 2) \norm{D}_2 + (2\alpha+1) \norm{u_2}_{L^\infty}) t \bigr)}, \quad t\in [0,T].
\end{equation*}
Using this estimate, and plugging it also in \eqref{eq:est_delta_m_dot}, we~get
\[ \| \delta m \|_{W^{1,\infty}} \le C(T, \| u_2 \|_{L^\infty}) \| \delta u \|_{L^\infty}, \]
and the conclusion follows.
\end{proof}

\begin{lemma} \label{lem:reach_fin_time}
If $\gamma_1 < \gamma_2$, there exists $\delta > 0$ such that for all $U > 0$, if $\abs{m (0) + e_1} < \delta$, then $- e_1$ can be reached in finite time with a control $u$ such that $\abs{u} \leq U$.
\end{lemma}

\begin{proof}
Let us introduce $m$ as the solution to \eqref{LL:ODE} with the feedback control term
\begin{equation*}
u (t) = \frac{U}{\sqrt{1 + \alpha^2}}\, \frac{\alpha (- e_1 + (e_1 \cdot m(t)) m(t)) + e_1 \wedge m(t)}{\sqrt{1- (m(t) \cdot e_1)^2}},
\end{equation*}
so that the equation on $m$ becomes autonomous, and is well defined as long as \hbox{$m (t) \neq \pm e_1$}. Observing that
\[
\biggl\{m,\frac{e_1-(m\cdot e_1)m}{\sqrt{1-(m\cdot e_1)^2}},\frac{m\wedge e_1}{\sqrt{1-(m\cdot e_1)^2}}\biggr\}
\]
is an orthonormal basis, one immediately gets that $|u(t)|=U$ for a.e. $t\in [0,T]$.

Let $m=(m_1,m_2,m_3)$ be the coordinates of $m$. From \eqref{LL:ODE}, the ODEs satisfied by~$m_2$ and~$m_3$ are
\begin{align*} \notag
\dot m_2 &= - \alpha [ (\gamma_2 - \gamma_1) m_2 - ((\gamma_2 - \gamma_1) m_2^2 + (\gamma_3 - \gamma_1) m_3^2) m_2 ] + (\gamma_1 - \gamma_3) m_1 m_3 + v_2, \\
\dot m_3 &= - \alpha [ (\gamma_3 - \gamma_1) m_3 - ((\gamma_2 - \gamma_1) m_2^2 + (\gamma_3 - \gamma_1) m_3^2) m_3 ] - (\gamma_1 - \gamma_2) m_1 m_2 + v_3.
\end{align*}
Therefore, by~setting $\widetilde{m} \coloneqq (m_2, m_3)$, it follows that $\widetilde{m}$ solves the controlled system
\begin{equation} \label{eq:ode_tilde_m_2}
\dot{\tilde{m}} = A_- \tilde{m} + \xi_- + \tilde v,
\end{equation}
where
\begin{align*}
A_- &=
\begin{bmatrix}
- \alpha (\gamma_2 - \gamma_1) & (\gamma_3 - \gamma_1) \\
- (\gamma_2 - \gamma_1) & - \alpha (\gamma_3 - \gamma_1)
\end{bmatrix}, \\
\xi_- &=
\begin{bmatrix}
\alpha ((\gamma_2 - \gamma_1) m_2^2 + (\gamma_3 - \gamma_1) m_3^2) m_2 - (\gamma_3 - \gamma_1) (1 + m_1) m_3 \\
\alpha ((\gamma_2 - \gamma_1) m_2^2 + (\gamma_3 - \gamma_1) m_3^2) m_3 + (\gamma_2 - \gamma_1) (1 + m_1) m_2
\end{bmatrix},
\end{align*}
and $\tilde v = (v_2, v_3)$ where $v = (v_1,v_2,v_3)=\alpha (u - (u \cdot m) m) - m \wedge u$, which means here
\begin{equation*}
v = U \sqrt{1 + \alpha^2}\,\frac{(- e_1 + m_1 (t) m(t))}{\sqrt{\abs{m(t)}^2 - (m_1 (t))^2}} = U \sqrt{1 + \alpha^2}\, \frac{(- e_1 + m_1 (t) m(t))}{\abs{\tilde{m}}},
\end{equation*}
We infer that $\tilde{v} = U \sqrt{1 + \alpha^2} m_1 \tilde{m} /\abs{\tilde{m}}$.
Observing that $ (1 - m_1) (1 + m_1) = \abs{\tilde{m}}^2$ yields, as soon as $m_1 \leq 0$, {$\abs{1 + m_1} \leq \abs{\tilde{m}}^2$ and thus}
\begin{equation*}
\abs{\xi_- (t)} \leq (1 + \abs{\alpha}) \, \delta \gamma_+ \, \abs{\tilde{m} (t)}^3,
\end{equation*}
where $\delta \gamma_+ \coloneqq \gamma_3 - \gamma_1 > 0$, and also
\begin{equation*}
\Babs{\tilde{v} + U \sqrt{1 + \alpha^2} \frac{\tilde{m} (t)}{\abs{\tilde{m}}}} \leq U \sqrt{1 + \alpha^2} \abs{\tilde{m} (t)}^2.
\end{equation*}
With these estimates and by taking the inner product of \eqref{eq:ode_tilde_m_2} with $\tilde{m}$, we~get
\begin{multline*}
\frac{1}{2} \frac{\diff}{\diff t} \abs{\tilde m (t)}^2 \leq - U \sqrt{1 + \alpha^2} \abs{\tilde{m} (t)} + \norm{A_-}_2 \abs{\tilde{m} (t)}^2 + U \sqrt{1 + \alpha^2} \abs{\tilde{m} (t)}^3\\
+ (1 + \abs{\alpha}) \, \delta \gamma_+ \, \abs{\tilde{m} (t)}^4,
\end{multline*}
and
\begin{equation*}
\frac{\diff}{\diff t} \abs{\tilde m (t)} \leq - U \sqrt{1 + \alpha^2} + \norm{A_-}_2 \abs{\tilde{m} (t)} + U \sqrt{1 + \alpha^2} \abs{\tilde{m} (t)}^2 + (1 + \abs{\alpha}) \, \delta \gamma_+ \, \abs{\tilde{m} (t)}^3.
\end{equation*}

Let us introduce $\delta_U \in (0,1/2)$ small enough (depending on $U > 0$) so that
\begin{equation}\label{choice:deltaU}
(U \sqrt{1 + \alpha^2} + \norm{A_-}_2) \delta_U + (1 + \abs{\alpha}) \, \delta \gamma_+ \, \delta_U^3 \leq \frac{U \sqrt{1 + \alpha^2}}{2}.
\end{equation}
Then, if $\abs{m (0) + e_1} < \delta_U$, which gives $\abs{\tilde{m} (0)} < \delta_U$, one has
\begin{equation*}
\frac{\diff}{\diff t} \abs{\tilde m (t)} \leq - \frac{U \sqrt{1 + \alpha^2}}{2}<0,
\end{equation*}
as long as $\abs{\tilde{m} (t)} < \delta_U$. This yields that, for such time intervals, the mapping $t\mto \abs{\tilde m (t)}$ is decreasing.

Therefore, this shows that if $\delta_U$ satisfies \eqref{choice:deltaU} and $m(0)$ is such that $|m(0)+e_1|<\delta_U$, then $\abs{\tilde{m} (t)} < \delta_U$ for all $t \geq 0$ and that $\tilde m (t)$ reaches $0$ in finite time.
In other words, $- e_1$ can be reached in finite time with a control $u$ such that $\abs{u} \leq U$ if $m$ is such that $\abs{m + e_1} < \delta_U$.

To conclude, it remains to drop the dependency of $\delta$ in $U$. Let us use that
$- e_1$ is asymptotically stable according to Proposition \ref{prop:asympt_stab}.
Therefore, there exists $\delta>0$ such that, starting from a point $m(0)$ chosen so that $\abs{m (0) + e_1} < \delta$, we~can first let the system evolve without control until we obtain $\abs{m (T_U) + e_1} < \delta_U$ for some finite time~$T_U$. From this moment, we~know we can reach $- e_1$ in finite time, whence the expected conclusion.
\end{proof}

Recall for the sake of readability that the notation $\mathcal{T}_U$ has been introduced in Section~\ref{sec:OCP}.
\begin{lemma} \label{lem:U_minus_eps}
Let $\gamma_1 < \gamma_2$ and $U > 0$ such that $\mathcal{T}_U < \infty$. Then there exists $\varepsilon > 0$ such that $\mathcal{T}_{U-\varepsilon} < \infty$.
\end{lemma}

\begin{proof}
Since $\mathcal{T}_U < \infty$ and according to Theorem~\ref{th:ex_min}, there exists $u_* \in L^\infty (0, \mathcal{T}_U)$ such that $m_* (0) = e_1$ and $m_* (\mathcal{T}_U) = - e_1$. Now, let us consider $m$ the solution to \eqref{LL:ODE} associated to the control choice $u = \frac{U - \varepsilon}{U} u_*$ for some $\varepsilon \in (0, U)$ to be defined later. From Lemma \ref{lem:cont_flow}, we~obtain
\begin{equation*}
\norm{m - m_*}_{W^{1, \infty} (0, \mathcal{T}_U)} \leq C \Bnorm{\frac{U - \varepsilon}{U} u_* - u_*}_{L^\infty (0, \mathcal{T}_U)} = C \varepsilon \frac{\norm{u_*}_{L^\infty (0, \mathcal{T}_U)}}{U} \leq C \varepsilon,
\end{equation*}
for some $C > 0$.

Since $m_* (\mathcal{T}_U) = - e_1$ by definition, we~can take $\varepsilon > 0$ small enough so that $\abs{m (\mathcal{T}_U) + e_1} < \delta$, where $\delta > 0$ is given by Lemma \ref{lem:reach_fin_time}. From this lemma, we~know we can reach $- e_1$ in finite time, and since $\abs{u} \leq U - \varepsilon$, this leads to $\mathcal{T}_{U-\varepsilon} < \infty$.
\end{proof}

\begin{lemma} \label{lem:T_U_decreas}
$\mathcal{T}_U$ is non-increasing with respect to $U > 0$. In particular, if $\mathcal{T}_{U_0} < \infty$ for some $U_0 > 0$, then $\mathcal{T}_U < \infty$ for all $U > U_0$.
\end{lemma}

\begin{proof}
This property is an immediate consequence of the definition of $\mathcal{T}_U$ and the fact that the sets $\mathcal{O}_{U }$ are increasing for the inclusion.
\end{proof}
\subsection{Emergence of a threshold}
The following result is the most crucial for concluding. It~quantifies the asymptotic stability of $e_1$ for the evolution of the magnetization $m$, with respect to $u$ viewed as a perturbation. It~is notable that its proof not only highlights the emergence of a threshold but also provides an explicit expression.
\begin{lemma} \label{lem:U_stab}
Let us assume that $\gamma_1 < \gamma_2$. There exists $U_{\textnormal{stab}}>0$ depending only on $\gamma_3-\gamma_1$, $\gamma_2-\gamma_1$ and $\alpha$ such that, for any $U < U_{\textnormal{stab}}$, the following holds. Let $(m,u)$ be a solution of \eqref{LL:ODE} on $[0,+\infty)$, such that $u \in L^\infty ([0, \infty))$ and $\norm{u}_{L^\infty} \leq U $. Then for all $t \ge 0$, $m_1(t) \ge 0$.
\end{lemma}

In other words, $m$ remains in the hemisphere with pole $e_1$, and in particular, $m$~can not reach $-e_1$. One has the same statement if $(m,u)$ are defined on a bounded interval $[0,T]$.

\begin{proof}
Let $v = \alpha (u - (u \cdot m) m) - m \wedge u$. Then, using that $v$ reads as the sum of two orthogonal terms, one has $\abs{v}^2 \leq (1 + \alpha^2) U ^2$.
Moreover, since $\abs{m}^2 = 1$, there holds
\begin{equation*}
Dm \cdot m = \gamma_1 + (\gamma_2 - \gamma_1) m_2^2 + (\gamma_3 - \gamma_1) m_3^2.
\end{equation*}
As in the proof of Lemma \ref{lem:reach_fin_time}, $\widetilde{m}$ solves the controlled system
\begin{equation} \label{eq:ode_tilde_m}
\dot{\widetilde{m}} = A \widetilde{m} + \xi + \tilde v.
\end{equation}
where
\begin{align} \label{def:A_xi}
A &=
\begin{bmatrix}
- \alpha (\gamma_2 - \gamma_1) & - (\gamma_3 - \gamma_1) \\
\gamma_2 - \gamma_1 & - \alpha (\gamma_3 - \gamma_1)
\end{bmatrix}, \\
\xi &=
\begin{bmatrix}
\alpha ((\gamma_2 - \gamma_1) m_2^2 + (\gamma_3 - \gamma_1) m_3^2) m_2 + (\gamma_3 - \gamma_1) (1 - m_1) m_3 \\
\alpha ((\gamma_2 - \gamma_1) m_2^2 + (\gamma_3 - \gamma_1) m_3^2) m_3 - (\gamma_2 - \gamma_1) (1 - m_1) m_2
\end{bmatrix}.\nonumber
\end{align}
It is important to note the sign change between $A$ and $\xi$ in this context, as well as~$A_-$ and~$\xi_-$ introduced in the proof of Lemma \ref{lem:reach_fin_time}.
Let $\nu \in (0, 1]$ to be fixed later and define
\[
T_\nu = \inf \{ t \geq 0 \ |\ |\widetilde{m}(t)| \geq \nu \}.
\]
Our goal is to derive suitable bounds on $\tilde m$, so that for a well chosen $\nu$, $T_\nu = +\infty$.

Since $m(0) = e_1$ and $m$ is continuous, we~know that $T_\nu > 0$. Note that one has necessarily $m_1(\cdot)>0$ on $(0,T_\nu)$. Then, for all $t \in [0, T_\nu)$, using that $m$ is normalized, there holds like previously
$
0 \leq 1 - m_1 (t)\leq 1 - m_1 (t)^2 = \abs{\widetilde{m} (t)}^2,
$
and therefore
\begin{equation*}
\abs{\xi (t)} \leq (1 + \abs{\alpha}) \, \delta \gamma_+ \, \abs{\widetilde{m} (t)}^3 \leq (1 + \abs{\alpha}) \, \delta \gamma_+ \, \nu^3,
\end{equation*}
where $\delta \gamma_+ \coloneqq \gamma_3 - \gamma_1 \geq \gamma_2 - \gamma_1 \eqqcolon \delta \gamma_- > 0$.
On the other hand, thanks to the Duhamel formula on \eqref{eq:ode_tilde_m} using the fact that $\widetilde{m} (0) = 0$, there holds
\begin{equation*}
\widetilde{m} (t) = \int_0^t \exp{\bigl((t-s) A \bigr)}\, (\xi (s) + \tilde v(s)) \diff s
\end{equation*}
for all $t \geq 0$. This, together with the previous estimates, drives to
\begin{equation}\label{m:1903}
\begin{aligned}
\abs{\widetilde{m} (t)} &\leq \int_0^t \norm{\exp{\bigl((t-s) A \bigr)}}_2 \Bigl((1 + \abs{\alpha}) \delta \gamma_+ \nu^3 + \sqrt{1 + \alpha^2} U \Bigr) \diff s \\
&\leq (1 + \abs{\alpha}) (\delta \gamma_+ \nu^3 + U) \int_0^t \norm{\exp{\bigl((t-s) A \bigr)}}_2 \diff s \\
&\leq (1 + \abs{\alpha}) (\delta \gamma_+ \nu^3 + U) \int_0^{t} \norm{\exp{\left(s A \right)}}_2 \diff s,
\end{aligned}
\end{equation}
for all $t \in [0, T_\nu)$, where $\norm{\cdot}_2$ still stands for the operator norm associated to the euclidean norm $\abs{\cdot}$.

We will now provide an estimate of the norm of the exponential matrix. Recall that the characteristic polynomial of $A$ is $P_A (X) = X^2 - \Tr (A) X + \det(A)$ with
\begin{equation*}
\det(A) = (1 + \alpha^2) \delta \gamma_- \, \delta \gamma_+ > 0, \qquad
\Tr (A) = - \alpha (\delta \gamma_- + \delta \gamma_+) < 0.
\end{equation*}
Its discriminant $\Delta$ reads $\Delta = \Tr (A)^2 - 4 \det(A) = \alpha^2 (\delta \gamma_+ - \delta \gamma_-)^2 - 4\delta \gamma_- \delta \gamma_+$. To~compute the eigenvalues of $A$, we~have to distinguish between several cases.

\subsubsection*{First case: $\Delta > 0$.}
Then its eigenvalues are $\lambda_\pm \coloneqq \frac12(\Tr (A) \pm \sqrt{\Delta})$.
Remark that both eigenvalues of $A$ are negative (according to the signs of the trace and the determinant above) and different from each other, which means that $A$ is diagonalizable. Therefore, we~infer\footnote{Here, the Lagrange interpolation formula is used to compute the exponential of $A$: for every matrix $M\in \mathcal{M}_d(\C)$ whose spectrum $\{\lambda_i\}_{1\leq i\leq d}$ consists of distinct eigenvalues, one has
\[
\exp(M)=\sum_{j=1}^de^{\lambda_j}\prod_{i\neq j}\frac{M-\lambda_i\operatorname{I}_d}{\lambda_j-\lambda_i} .
\]
} that
\begin{align*}
\exp{(sA)} &= e^{s \lambda_+} \frac{sA - s \lambda_- \operatorname{I}_2}{s \lambda_+ - s \lambda_-} + e^{s \lambda_-} \frac{sA - s \lambda_+ \operatorname{I}_2}{s \lambda_- - s \lambda_+} \\
&= \frac{1}{\sqrt{\Delta}} \left(e^{s \lambda_+} (A - \lambda_- \operatorname{I}_2) - e^{s \lambda_-} (A - \lambda_+ \operatorname{I}_2) \right) \\
&= \frac{1}{\sqrt{\Delta}} \left((e^{s \lambda_+} - e^{s \lambda_-}) A + (e^{s \lambda_-} \lambda_+ - e^{s \lambda_+} \lambda_-) \operatorname{I}_2 \right) \\
&= \frac{e^{s \lambda_+}}{\sqrt{\Delta}} \left((1 - e^{- s \sqrt{\Delta}}) A + (e^{- s \sqrt{\Delta}} \lambda_+ - \lambda_-) \operatorname{I}_2 \right) \\
&= s e^{s \lambda_+} \Bigl(\frac{1 - e^{- s \sqrt{\Delta}}}{s \sqrt{\Delta}} A - \lambda_- \frac{1 - e^{- s \sqrt{\Delta}}}{s \sqrt{\Delta}} \operatorname{I}_2 \Bigr) + e^{s \lambda_-} \operatorname{I}_2.
\end{align*}
Thus, using the facts that $\lambda_- < \lambda_+ < 0$, $\norm{A}_2 \geq \abs{\lambda_-}$ and also that the function~$f$ given by $f(x) = \psfrac{1 - e^{-x}}{x}$ analytically extended to $\R$ is uniformly bounded by~$1$ on $[0, \infty)$, we~get
\begin{equation*}
\norm{\exp{(sA)}}_2 \leq e^{s \lambda_+} \bigl(s (\norm{A}_2 + \abs{\lambda_-}) + 1 \bigr) \leq e^{s \lambda_+} \left(2 s \norm{A}_2 + 1 \right).
\end{equation*}
Hence,
\begin{align*}
\int_0^{t} \norm{\exp{\left(s A \right)}}_2 \diff s &\leq \abs{\lambda_+}^{-1} (1 - e^{\lambda_+ t}) + 2 \norm{A}_2 \abs{\lambda_+}^{-2} (1 - (|\lambda_+| t + 1) e^{\lambda_+ t}) \\
&\leq (1 - e^{\lambda_+ t}) \left(\abs{\lambda_+}^{-1} + 2 \norm{A}_2 \abs{\lambda_+}^{-2} \right) \\
&\leq 3 (1 - e^{\lambda_+ t}) \norm{A}_2 \abs{\lambda_+}^{-2},
\end{align*}
and according to \eqref{m:1903}, one has for all $t \in [0, T_\nu)$
\begin{equation*}
\abs{\widetilde{m} (t)} \leq 3 \norm{A}_2 \abs{\lambda_+}^{-2} (1 - e^{\lambda_+ t}) (1 + \abs{\alpha}) (\delta \gamma_+ \nu^3 + U).
\end{equation*}
To conclude, we~will choose $U $ adequately so that the function
\[
x \mto 3 \norm{A}_2 \abs{\lambda_+}^{-2} (1 + \abs{\alpha}) (\delta \gamma_+ x^3 + U)
\]
admits a fixed point $x_0$ in $(0, 1]$.
This is possible thanks to the next lemma, whose proof is postponed to the end of this section for the sake of clarity.
\begin{lemma} \label{lem:fixed_point}
Let $a, b, c > 0$. The function $x \mto a^{-1} (b x^3 + c)$ has a fixed point $x_0$ in $(0, 1]$ if, and only if $c \leq a x_1 - b x_1^3$ where $x_1 = \min \{1, \sqrt{\sfrac{a}{3 b}} \}$.
\end{lemma}

Remark that if $x_1$ is as in this lemma, one has $a x_1 - b x_1^3 \geq \frac{2}{3} a x_1>0$.
{Thus, by~applying Lemma~\ref{lem:fixed_point} with the parameters $a=|\lambda_+|^2/(3\norm{A}_2(1+|\alpha|))$, $b=\delta \gamma_+$ and $c=U $, we~arrive at the assumption that
\begin{equation*}
U \leq \frac{\abs{\lambda_+}^2}{3 \norm{A}_2 (1 + \abs{\alpha})} x_1 - \delta \gamma_+ \, x_1^3, \quad\text{with}\quad x_1 \coloneqq \min \biggl\{ 1, \frac{\abs{\lambda_+}}{3 \sqrt{\norm{A}_2 (1 + \abs{\alpha}) \delta \gamma_+}} \biggr\}.
\end{equation*}
Therefore, based on the conclusion of this lemma, we~can set $\nu = x_0$, and the preceding estimate yields
\begin{equation*}
\abs{\widetilde{m}(t)} \leq (1 - e^{\lambda_+ t}) \nu,
\end{equation*}
for all $t \in [0, T_\nu)$.} A continuity argument then implies that $T_\nu = \infty$. In other words, for all $t \geq 0$,
\[
|m_1(t)| = \sqrt{1-|\tilde m(t)|^2} \ge \sqrt{1-(1 - e^{\lambda_+ t})^2 \nu^2} >0.
\]
Now, $m_1$ is continuous, so that it keeps a constant sign. As $m_1(0)=1$, $m_1(t) \ge 0$ for all $t \ge 0$, which is the desired conclusion.

\subsubsection*{Second case: $\Delta < 0$.}
In this case, the eigenvalues are
\begin{equation*}
\lambda_\pm \coloneqq \frac{\Tr (A) \pm i \sqrt{- \Delta}}{2}.
\end{equation*}
One more time, the two eigenvalues are distinct, complex conjugate with negative real part. Yet, the same decompositions as previously can still be applied and there holds
\begin{align*}
\exp(s&A) = e^{s \lambda_+} \frac{sA - s \lambda_- \operatorname{I}_2}{s \lambda_+ - s \lambda_-} + e^{s \lambda_-} \frac{sA - s \lambda_+ \operatorname{I}_2}{s \lambda_- - s \lambda_+} \\
&= \frac{1}{i \sqrt{- \Delta}} \left(e^{s \lambda_+} (A - \lambda_- \operatorname{I}_2) - e^{s \lambda_-} (A - \lambda_+ \operatorname{I}_2) \right) \\
&= \frac{1}{i \sqrt{- \Delta}} \left((e^{s \lambda_+} - e^{s \lambda_-}) A + (e^{s \lambda_-} \lambda_+ - e^{s \lambda_+} \lambda_-) \operatorname{I}_2 \right) \\
&= \frac{e^{\Psfrac{s}{2} \Tr (A)}}{i \sqrt{- \Delta}} \left((e^{\sfrac{i s \sqrt{- \Delta}}{2}} - e^{- \sfrac{i s \sqrt{- \Delta}}{2}}) A + (\lambda_+ e^{- \sfrac{i s \sqrt{- \Delta}}{2}} - \lambda_- e^{\sfrac{i s \sqrt{- \Delta}}{2}}) \operatorname{I}_2 \right) \\
&= \frac{s}{2} e^{\Psfrac{s}{2} \Tr (A)} \Bigl[ 2 A - \Tr (A) \operatorname{I}_2 \Bigr] \sinc{\Bigl(\frac{s \sqrt{- \Delta}}{2} \Bigr)} + e^{\Psfrac{s}{2} \Tr (A)} \cos{\Bigl(\frac{s \sqrt{- \Delta}}{2} \Bigr)} \operatorname{I}_2.
\end{align*}
Thus, we~get
\begin{equation*}
\norm{\exp{(sA)}}_2 \leq \frac{s}{2} e^{\Psfrac{s}{2} \Tr (A)} \left(2 \norm{A}_2 - \Tr (A) \right) + e^{\Psfrac{s}{2} \Tr (A)} \leq e^{\Psfrac{s}{2} \Tr (A)} \left(2 s \norm{A}_2 + 1 \right),
\end{equation*}
since $\Tr (A) = \lambda_+ + \lambda_-$ and $\abs{\lambda_\pm} \leq \norm{A}_2$.
Hence, following the same way as in the first case, we~get
\begin{align*}
\int_0^{t} \norm{\exp{\left(s A \right)}}_2 \diff s &\leq (1 - e^{\Psfrac{1}{2} \Tr (A) t}) \left(2 \abs{\Tr (A)}^{-1} + 8 \norm{A}_2 \abs{\Tr (A)}^{-2} \right) \\
&\leq 12 (1 - e^{\Psfrac{1}{2} \Tr (A) t}) \norm{A}_2 \abs{\Tr (A)}^{-2},
\end{align*}
and according to \eqref{m:1903}, one has for all $t \in [0, T_\nu)$
\begin{equation*}
\abs{\widetilde{m} (t)} \leq 12 \norm{A}_2 \abs{\Tr (A)}^{-2} (1 - e^{\Psfrac{1}{2}\Tr (A) t}) (1 + \abs{\alpha}) (\delta \gamma_+ \nu^3 + U).
\end{equation*}
Now, by~mimicking the reasoning done in the first case, by~assuming
\begin{equation*}
U \leq \frac{\Tr (A)^2}{12 \norm{A}_2 (1 + \abs{\alpha})}\, x_1 - \delta \gamma_+ \, x_1^3, \quad\text{with}\quad x_1 \coloneqq \min \biggl\{ 1, \frac{\Tr (A)}{6 \sqrt{\norm{A}_2 (1 + \abs{\alpha}) \delta \gamma_+}} \biggr\},
\end{equation*}
and taking $\nu = x_0$ given by Lemma~\ref{lem:fixed_point}, the previous estimate leads to
\begin{equation*}
\abs{\widetilde{m}(t)} \leq (1 - e^{\Psfrac{1}{2} \Tr (A) t}) \nu,
\end{equation*}
for all $t \in [0, T_\nu)$. Arguing as in the first case, we~infer that $T_\nu = \infty$ in this case as well, and then, $m_1(t) > 0$ for all $t \ge 0$.

\subsubsection*{Third case: $\Delta = 0$.}
In this case, both eigenvalues are equal, one has \hbox{$\lambda \!=\! \Tr (A)/2\!<\!0$}. Note that, in~that case, $A - \frac{1}{2} \Tr (A) \operatorname{I}_2$ is therefore a non-zero nilpotent matrix, and more precisely $(A - \frac{1}{2}\Tr (A) \operatorname{I}_2)^2 = (s A - \frac{s}{2} \Tr (A) \operatorname{I}_2)^2 = 0$. Thus, there holds
\begin{align*}
\exp{(sA)} &= \exp{\left(\tfrac{s}{2} \Tr (A) \operatorname{I}_2 \right)} \, \exp{(sA - \tfrac{s}{2} \Tr (A) \operatorname{I}_2)} \\
&= e^{\Psfrac{s}{2} \Tr (A)} (\operatorname{I}_2 + sA - \tfrac{s}{2} \Tr (A) \operatorname{I}_2),
\end{align*}
which yields
\begin{equation*}
\norm{\exp{(sA)}}_2 \leq e^{\Psfrac{s}{2} \Tr (A)} (1 + s (\norm{A}_2 - \tfrac{1}{2} \Tr (A))) \leq e^{\Psfrac{s}{2} \Tr (A)} (1 + 2 s \norm{A}_2).
\end{equation*}
The computations are then exactly the same ones as in the second case, and the conclusion follows in the same fashion.
\end{proof}

\begin{proof}[Proof of Lemma \ref{lem:fixed_point}]
We are looking for a root $x_0 \in (0, 1]$ of the polynomial function~$f$ given by $f (X) = b X^3 - a X + c$, whose derivative $3 b X^2 - a$ is negative for $X < \sqrt{\sfrac{a}{3 b}} \eqqcolon x_2$ and positive for $X > x_2$. The minimum in $[0, 1]$ is therefore reached at $x_1$ and, since $f (0) = c > 0$, there is a root if and only if $f(x_1) \leq 0$, which corresponds to the assumption in the statement.
\end{proof}

\begin{proof}[Proof of Theorem \ref{th:ex_U_crit}]

Define
\begin{equation*}
U_\textnormal{crit} \coloneqq \inf \{ U \, | \, \mathcal{T}_U < \infty \}.
\end{equation*}
From Lemma \ref{lem:U_stab}, we~know that $U_\textnormal{crit} > 0$. Lemma~\ref{lem:T_U_decreas} proves the second point. To~investigate the case where $U =U_\textnormal{crit}$, observe that, by~definition, \eqref{OCP:minT} has no solution for all $U \in (0, U_\textnormal{crit})$. Moreover, if \eqref{OCP:minT} had a solution for $U = U_\textnormal{crit}$, then Lemma~\ref{lem:U_minus_eps} would provide a contradiction with respect to the definition of $U_\textnormal{crit}$.
\end{proof}
\appendix
\setcounter{section}{1}
\setcounter{theorem}{24}
\section{Stability of steady-states for Equation~\texorpdfstring{\eqref{LL:ODE}}{(3)}}\label{append:e1AS}
Let us first notice that, if $\bar m$ is a steady-state of Equation~\eqref{LL:ODE} and if no control is applied on this system ($h_{\rm ext}=0$), then by orthogonality of the terms in the right-hand side, it satisfies
\[
h_0(\bar m)=(h_0(\bar m)\cdot \bar m)\bar m\quad \text{and}\quad \bar m\wedge h_0(\bar m)=0.
\]
Therefore, $\bar m$ is an eigenfunction of $D\bar m$ and we infer that $\bar m=\pm e_j$, $j=1,2,3$ whenever $\gamma_1<\gamma_2\leq \gamma_3$.
\begin{proposition}[Asymptotic stability] \label{prop:asympt_stab}
Let $\gamma_1, \gamma_2, \gamma_3$ be sorted in ascending order $\gamma_1<\gamma_2\leq\gamma_3$. Then, in~the absence of any control $h_{\rm ext}$, the steady-state $\pm e_1$ is an asymptotically stable equilibrium state for Equation~\eqref{LL:ODE}. Nevertheless, the steady-states $\pm e_2$ and $\pm e_3$ are linearly unstable steady-states for Equation~\eqref{LL:ODE}.
\end{proposition}
\begin{proof}
Let $h=(h_1, h_2, h_3)^T\in\R^3$ be a small perturbation such that $e_1+h$ is still an admissible magnetization, \ie on the unit sphere $\mathbb{S}^2\subset\R^3$. We obtain:
\begin{equation*}
\|e_1+h\|^2=1\iff (1+h_1)^2+h_2^2+h_3^2=1
\iff h_1=-\frac{1}{2}\left(h_1^2+h_2^2+h_3^3\right)=\mathcal{O}(\|h\|^2).
\end{equation*}
The unknown $h_1$ is therefore of second order and does not occur in a linearized system of the first order.

By linearizing Equation~\eqref{LL:ODE} around the equilibrium state $e_1$, one has, without any control $u$:
\begin{equation}
\left\{
\begin{aligned}
\dot{h_2} & =\alpha(\gamma_1-\gamma_2)h_2+(\gamma_1-\gamma_3)h_3+\operatorname{O}(\|h\|^2),\\
\dot{h_3} & =\alpha(\gamma_1-\gamma_3)h_3+(\gamma_2-\gamma_1)h_2+\operatorname{O}(\|h\|^2).\\
\end{aligned}
\right.
\end{equation}

The Jacobian matrix of the linearized system around $e_1$ is therefore :
\[ J=\begin{pmatrix}\alpha(\gamma_1-\gamma_2)&\gamma_1-\gamma_3\\
\gamma_2-\gamma_1&\alpha(\gamma_1-\gamma_3)\end{pmatrix}.\]

Since $\gamma_1<\gamma_2\leq\gamma_3$, one has
\[
\det (J)\!=\!(\alpha^2+1)(\gamma_1-\gamma_2)(\gamma_1-\gamma_3)\!>\!0\quad \text{and}\quad
\operatorname{Tr}(J)\!=\!\alpha\bigl[(\gamma_1-\gamma_2)+(\gamma_1-\gamma_3)\bigr]\!<\!0.
\]
We infer that the two eigenvalues of the Jacobian matrix are of negative real parts. The steady state $e_1$ is therefore linearly stable and is {a hyperbolic point} (no eigenvalue with zero real part).

{The Hartman-Grobman theorem} \cite{perko2013differential} allows to conclude about the asymptotic stability of $e_1$ for the non-linear Equation~\eqref{LL:ODE} without any control $u$. As for $- e_1$, similar computations give the conclusion. Regarding now the stability of $\pm e_k$, $k=2,3$, notice that a similar computation drives to the following expression of the Jacobian determinant: $\det J=(\alpha^2+1)(\gamma_2-\gamma_1)(\gamma_2-\gamma_3)<0$. The expected conclusion follows.
\end{proof}

\begin{remark}
If $\gamma_1\leq\gamma_2\leq\gamma_3$, an eigenvalue of the Jacobian matrix may have a zero real part. In which case one can conclude that $e_1$ is linearly (non-asymptotically) stable, but {the Hartman-Grobman theorem} no longer applies to return to the non-linear Equation~\eqref{LL:ODE}.
\end{remark}
\begin{thebibliography}{BGR23}
\bibitem[15]{hubert2008magnetic} A. Hubert \& R. Schäfer – Magnetic domains: the analysis of magnetic microstructures, Springer Science \& Business Media, 2008.
\bibitem[7]{brown1963micromagnetics} W. F. Brown – Micromagnetics, Interscience, 1963.
\bibitem[19]{perko2013differential} L. Perko – Differential equations and dynamical systems, vol. 7, Springer Science \& Business Media, 2013.
\bibitem[3]{alouges2009magnetization} F. Alouges \& K. Beauchard – “Magnetization switching on small ferromagnetic ellipsoidal samples”, ESAIM Control Optim. Calc. Var. 15 (2009), no. 3, p. 676–711.
\end{thebibliography}
\end{document}
